% Verbatim copy of /Users/alingalatan/workspaces/grothe/handoff/manuscript/legacy/ch37-gpt-response-stability.tex (file dated 2026-10-25). Only this line was added.
% Consolidated source module (formerly Chapter 37). Owner: lower-bound/06-original-responses.tex
\section{Lower-bound GPT: response stability of near-maximizers, exact K for one-coordinate remote cubics, and complex orbit constraints (20 October)}\label{ch:response-stability}
\chapsummary{Near-maximizer constraints, remote-cubic special cases and complex orbit tools. Stationarity belongs to the original unconstrained norm problem in fixed dimension, not to every fixed-row response.}

\paragraph{Provenance and status (IDE note, 20 October).} Lower-bound GPT's three fragments \code{chapter\_fresh\_second\_chaos.tex}, \code{core\_fibre\_theorem.tex} and
\code{chapter\_complex\_orbits.tex} from \pathref{gpt-reviews/k-forest-20261020/}, ported verbatim apart from this header and the demotion of the orbit chapter to a section.
\emph{Our replay} \st{C}: \code{verify\_all.py} under \code{python -I} in an isolated venv with the pinned \code{python-flint 0.9.0} and \code{sympy 1.14.0}, from an unrelated directory:
\textbf{ALL SIX REPLAYS PASS in 6.8\,s}. These are exact identity and certificate checks (Arb profile $L^1$, peak anchor, core-fibre cone, fresh second-chaos coefficients and saddle,
2{,}929 Beta/Hermite cases and 8{,}379 singular-value duals, the weighted Bessel Schur identities), not a large cover. Log \pathref{gpt-reviews/k-forest-20261020/replay\_ide.log}.
The strategy notes (INTEGRATION.md) are in the package.
\providecommand{\eps}{\varepsilon}\providecommand{\ph}{\operatorname{phase}}

% Self-contained section fragment; requires amsmath and amssymb.
\subsection{A fresh second-chaos stability cut on actual response pairs}
\label{sec:fresh-second-chaos-stability}

This section proves a necessary condition on near-maximizers of the
original cubic witness. It does not prove Conjecture K or a numerical
Grothendieck endpoint. Put $T=P-\varepsilon Q$, with
$P=P_{1,0}$, $Q=P_{2,1}$ and $0<\varepsilon\le1$, and let
\[
 M=\sup_{n,\ |f|,|g|\le1}\operatorname{Re}\langle f,Tg\rangle.
\]
Scalar products are conjugate-linear in their first argument. The
supremum is over all finite-dimensional proper complex Gaussian spaces.

\paragraph{Exact finite-amplitude theorem.}
Adjoin an independent standard proper complex Gaussian $U$ and put
$Y=|U|^2-1$. For arbitrary real measurable multipliers $A,B$ on the
original coordinates, set
\[
 f_t=f e^{itAY},\qquad g_t=g e^{itBY},\qquad
 m_0(x)=\frac{e^{-ix}}{1-ix},\qquad
 m_1(x)=\frac{ix e^{-ix}}{(1-ix)^2}.
\]
Then, exactly,
\begin{align*}
 Pf_t&=P(fm_0(tA)),\\
 Qf_t&=Q(fm_0(tA))+Y P(fm_1(tA)).
\end{align*}
The projections on the right act in the original coordinates. Hence
\begin{equation}\label{eq:fresh-exact-value}
 J_t:=\operatorname{Re}\langle f_t,Tg_t\rangle
 =\operatorname{Re}\langle fm_0(tA),T(gm_0(tB))\rangle
 -\varepsilon\operatorname{Re}\langle P(fm_1(tA)),P(gm_1(tB))\rangle.
\end{equation}
In particular, $J_t\le M$. If $J_0\ge M-\delta$, this gives the finite cut
\begin{equation}\label{eq:fresh-finite-cut}
 \frac{J_t+J_{-t}}2-J_0\le\delta.
\end{equation}
There is no attainment assumption, limiting Gaussian replacement or
tail truncation in these statements.

For the proof, all Hermite modes of a radial function of $U$ have
bidegree $(j,j)$. Only $j=0,1$ can contribute to $P$ or $Q$;
their normalized modes are $1$ and $Y$. Direct integration gives
$\mathbb E e^{ixY}=m_0(x)$ and
$\mathbb E[Y e^{ixY}]=m_1(x)$. Product Hermite orthogonality proves the
two projection identities, and the two terms of $Qf_t$ are orthogonal.
Also $|m_0|\le1$ and $|m_1|\le1/2$, so the formula remains integrable
for unbounded measurable multipliers. No charge-homogeneity assumption
on $A,B,f,g$ is needed.

These cuts give a direct upper-bound dual on the original norm.
Let $C_j=(J_{t_j}+J_{-t_j})/2-J_0$ be finitely many such gains and
let $\lambda_j\ge0$ have a uniformly bounded sum. The tests and prices
may depend on the response pair. Then
\[
 M\le\sup_{f,g,n}\left\{J_0-\sum_j\lambda_j C_j\right\}.
\]
Indeed, along any sequence $J_0\to M$, each $C_j\le M-J_0$,
so the expression in braces has lower limit at least $M$.
Thus a global certificate for this bracket would bound the norm
directly, without requiring Conjecture K for every arbitrary residual.
Keeping only positive gains sharpens this to the exact identity
\[
 \boxed{M=\sup_{f,g,n}\left\{J_0-
                     \sum_j\lambda_j(C_j)_+\right\}.}
\]
The right side is at most $M$ by nonnegativity of the penalty, while
the preceding near-maximizer argument proves the reverse inequality.
Pair-dependent prices require a uniformly bounded total. A maximum
of finitely many positive gains may be used instead of their sum.

\paragraph{Quantitative Hessian theorem.}
Suppose $|A|,|B|\le K$, and define the real weights and quadratic form
\[
 a=\operatorname{Re}(\overline f\,Tg),\quad
 b=\operatorname{Re}(\overline g\,Tf),\qquad
 H(A,B)=\mathbb E(aA^2+bB^2)
       +2\varepsilon\operatorname{Re}\langle P(fA),P(gB)\rangle.
\]
The weights need not be nonnegative. For every real $t$,
\begin{equation}\label{eq:fresh-fourth-remainder}
 \left|\frac{J_t+J_{-t}}2-J_0+\frac{t^2}2H(A,B)\right|
 \le\left(\frac{13}4+3\varepsilon\right)K^4t^4.
\end{equation}
Consequently every $\delta$-near-maximizer satisfies
\begin{equation}\label{eq:fresh-near-max-hessian}
 \boxed{H(A,B)\ge-4\sqrt{(13/4+3\varepsilon)\delta}\,K^2
                    \ge-10K^2\sqrt\delta.}
\end{equation}

Here is a global remainder proof using only
$\mathbb EY=0$, $\mathbb EY^2=1$ and $\mathbb EY^4=9$.
The symmetrized payoff is the sum of the cosine-pair and sine-pair
payoffs. In the cosine pair, the two single increments have total
fourth-order remainder at most $3K^4t^4/4$, because
$|\cos x-1+x^2/2|\le x^4/24$ and
$\|Tf\|_1,\|Tg\|_1\le1$. The product of the two cosine increments
has absolute value at most $9K^4t^4/4$, using
$|\cos x-1|\le x^2/2$ and $\|T\|_{2\to2}\le1$.

For the sine pair write
$s_0(A)=\mathbb E_Y\sin(tAY)$ and
$s_1(A)=\mathbb E_Y[Y\sin(tAY)]$. The elementary bounds are
\[
 |s_0(A)|\le t^2K^2/2,\qquad |s_1(A)|\le |t|K,
 \qquad |s_1(A)-tA|\le3|t|^3K^3/2.
\]
The old-coordinate sine contribution is therefore at most
$K^4t^4/4$ in absolute value. By the same projection decomposition as
above, the new-coordinate contribution is
$-\varepsilon\operatorname{Re}\langle P(fs_1(A)),P(gs_1(B))\rangle$;
its error from the quadratic term is at most
$3\varepsilon K^4t^4$. The four remainder constants sum to
$13/4+3\varepsilon$, proving \eqref{eq:fresh-fourth-remainder}.
If $H=-h<0$, maximize $ht^2/2-CK^4t^4$ in $t^2$. Its maximum is
$h^2/(16CK^4)$ and cannot exceed $\delta$ by
\eqref{eq:fresh-finite-cut}; this proves
\eqref{eq:fresh-near-max-hessian}.

The necessary inequality remains valid if its weights $a,b$ are
replaced by $|Tg|,|Tf|$, since this only enlarges $H$. Thus the
positive-weight form applies to every near-maximizer, even without
exact best-response identities. For real multipliers in $L^8$ the same
argument replaces $(13/4+3\varepsilon)K^4$ by
\begin{align*}
 C(A,B)={}&\frac38\{\mathbb E(|Tg|A^4)+\mathbb E(|Tf|B^4)\}
       +\frac52\|A\|_4^2\|B\|_4^2\\
 &+\frac{3\varepsilon}2
       (\|A\|_6^3\|B\|_2+\|A\|_2\|B\|_6^3).
\end{align*}
The two weighted expectations are at most $\|A\|_8^4$ and
$\|B\|_8^4$. This gives an explicit Gaussian-moment version for
polynomial multipliers, with $H\ge-4\sqrt{C(A,B)\delta}$.
Equivalently the explicit penalty
\[
 D(A,B)=\frac{H(A,B)_-^2}{16C(A,B)}\le M-J_0
\]
gives, for every fixed finite $\lambda\ge0$,
\[
 \boxed{M=\sup_{f,g,n}\{J_0-\lambda D(A,B)\}.}
\]
Set $D=0$ when $C=0$; then the relevant multiplier norms and $H$
vanish. Finite sums with nonnegative bounded-total prices and the
absolute-weight version of $H$ are permitted. This is an exact
reformulation of the original norm, not an added hypothesis about
arbitrary cubic perturbations.

For finite bounded real features $A_i,B_j$, set
\[
 D_f[i,j]=\mathbb E(|Tg|A_iA_j),\quad
 D_g[i,j]=\mathbb E(|Tf|B_iB_j),\quad
 C[i,j]=\operatorname{Re}\langle P(fA_i),P(gB_j)\rangle.
\]
If $\sum_iA_i^2\le K^2$ and $\sum_jB_j^2\le K^2$ pointwise, then
every $\delta$-near-maximizer satisfies the finite Gram constraint
\[
 \begin{pmatrix}D_f&\varepsilon C\\
                 \varepsilon C^{\mathsf T}&D_g\end{pmatrix}
 \succeq-10K^2\sqrt\delta\,I.
\]
One implementable family is the clipped weighted scores
$\operatorname{clip}(\operatorname{Re}(\overline f\ell_i)/
(|Tg|+\tau),-K_i,K_i)$, with $\tau>0$, and their $g$ analogues.

\paragraph{Weighted first-chaos matrix.}
At an attained dimension-free maximum, replace the responses by phase
best responses, so $f=\operatorname{phase}(Tg)$,
$g=\operatorname{phase}(Tf)$, $a=|Tg|$, $b=|Tf|$.
On the realification of the first-chaos space, define
\[
 M_f(u,v)=\mathbb E\frac{
 \operatorname{Re}(\overline f\,\ell_u)
 \operatorname{Re}(\overline f\,\ell_v)}{a},
 \qquad
 M_g(u,v)=\mathbb E\frac{
 \operatorname{Re}(\overline g\,\ell_u)
 \operatorname{Re}(\overline g\,\ell_v)}{b}.
\]
When the matrices are finite, a necessary condition is
\begin{equation}\label{eq:fresh-matrix}
 \boxed{\varepsilon^2\lambda_{\max}
       (M_f^{1/2}M_gM_f^{1/2})\le1.}
\end{equation}
Indeed the zero-error Hessian inequality, with scaling and a sign change
of $B$, gives
\[
 \varepsilon|\operatorname{Re}\langle P(fA),P(gB)\rangle|
 \le\sqrt{\mathbb EaA^2\,\mathbb EbB^2}.
\]
The adjoint of $A\mapsto P(fA)$ in real $L^2(a\,d\gamma)$ is
$u\mapsto\operatorname{Re}(\overline f\ell_u)/a$, whose covariance
matrix is $M_f$. Truncation of multipliers proves the assertion for
the weighted spaces. This operator formulation handles singular
matrices without writing inverses. For near-maximizers one must use
bounded multipliers and the error in
\eqref{eq:fresh-near-max-hessian}, rather than silently imposing
\eqref{eq:fresh-matrix}. A fixed-dimensional maximizer alone need not
satisfy \eqref{eq:fresh-matrix}, since adjoining $U$ may improve its value.

At the diffuse Gaussian critical response, a first-chaos direction
independent of the response has
$M_f=M_g=\varepsilon^{-1}I$, by the scalar buffer stationarity equation
$\mathbb E|L+\sqrt vW|^{-1}=2/\varepsilon$. Thus the new matrix constraint
is exactly active in unused remote directions. This contrasts with the
strict slack of radial weighted-Bessel and global orbit constraints.
It does not show that every shared non-Gaussian core violates the cut.

\paragraph{An exact exclusion test.}
Take $\varepsilon=1$ and $f=g=\operatorname{phase}(Z)$. Then
\[
 Pf=\frac{\sqrt\pi}2 Z,\qquad
 Qf=-\frac{\sqrt\pi}{4\sqrt2}H(Z),\qquad
 Tf=\frac{\sqrt\pi}2 Z\left(\frac12+\frac{|Z|^2}4\right).
\]
This pair is a mutual phase best response with value $7\pi/32$.
For the bounded multipliers $A=1$, $B=-1$, however,
$H(A,B)=-\pi/16<0$. More concretely, $t=1/8$ in
\eqref{eq:fresh-exact-value} gives $J_t-J_0>27/32000$.
The standalone checker verifies this by rational alternating-series
enclosures for $\cos(1/4)$ and $\sin(1/4)$ and the inequality $\pi>3$.
This proves that the cut supplies information beyond phase
best-response and first-order stationarity.

There is also an exact bounded exclusion at the actual peak price
$\varepsilon=1891/2500$. Add an unused remote first-chaos Gaussian $X$
and put
\[
 u(T)=\min(T^{-1/2},100),\qquad
 A=\operatorname{Re}(X\overline{\operatorname{phase}(Z)})
       u(|Z|^2)\,\mathbf1_{\{|X|\le4\}},\qquad B=-A.
\]
Here $|A|\le400$. With $j=1-17e^{-16}$ and $I=\mathbb Eu(T)$,
direct angular integration gives
\[
 H(A,-A)=j\mathbb E[|T\operatorname{phase}(Z)|u(T)^2]
                 -\frac\varepsilon2j^2I^2.
\]
The elementary estimates
$j\ge1-17/65536$, $I\ge\sqrt\pi-1/50>7/4$, and
\[
 \mathbb E[|T\operatorname{phase}(Z)|u(T)^2]
 \le\frac\pi2\left(1-\frac{3\varepsilon}8\right)
\]
together with $\pi<22/7$ give $H(A,-A)<-1/32$ by rational
arithmetic. The checker verifies the last comparison and the
needed rational bounds on $\pi$. The large clipping bound makes
this an exact separation example, not evidence that the coarse
near-maximizer error constant is efficient in a numerical certificate.

\paragraph{Scope.}
General nonlinear measure-preserving first variations offer no free
shortcut: wherever $Tg\ne0$, differentiation of
$f=\operatorname{phase}(Tg)$ gives
$\operatorname{Re}(\overline{Tg}\,Df)=0$ pointwise. Linear unitary
rotations are useful because they commute with $P,Q$ and rewrite this
identity in the penalized coefficients. Nonlinear transports introduce
unpriced higher-chaos commutators. The present theorem instead uses
simultaneous response perturbations and an exact new $(1,1)$ mode.

To obtain a numerical endpoint one still has to select bounded
multipliers detecting the dangerous shared-sector region, combine the
cut with the response-consistency dual, and certify a complete row
cover. No such global optimization is claimed here.

% Proposed manuscript insertion; this is an input fragment, not a replacement main file.
% Requires amsmath, amssymb. No project-specific macros are used.
\subsection{An exact remote comparison joining a coherent core and latent factors}
\label{sec:remote-core-factor-cone}

This result concerns a residual independent of the visible coordinate. It does
not settle the shared-visible $A/B$ obligation, and changes no global endpoint.
It does combine a coherent cubic core with arbitrary latent Gaussian factors
without a central-limit error.

Suppose $L\overset d=\sqrt V Z$, where $Z$ is a standard proper complex Gaussian,
$V\ge0$, and $\mathbb E\sqrt V<\infty$. Let $X,G$ be standard proper complex
Gaussians. Assume that $X$, $G$, $L$, and the jointly square-integrable
coefficient vector $(c,C,D,U,A,Q_0)$ are mutually independent random objects. Put
\[
 H(X)=\frac{X(|X|^2-2)}{\sqrt2},\qquad
 Q=cH(X)+CX+D\overline X+\frac{UX^2}{\sqrt2}
          +A(|X|^2-1)+Q_0.
\]
The coefficients may be arbitrarily dependent on each other. Orthogonality in
$X$ gives
\[
 v:=\mathbb E|Q|^2
 =\mathbb E\bigl(|c|^2+|C|^2+|D|^2+|U|^2+|A|^2+|Q_0|^2\bigr).
\]
For $g\ge0$ define
\[
 f(v)=\sqrt\pi\,\mathbb E\sqrt{V+g^2+v},\qquad
 d=-\frac v2f''(v),\quad
 d_3=-vf''(v)-\frac{v^2}2f'''(v),\quad
 d_0=\frac{f(v)-f(0)}v-f'(v).
\]
For $v>0$, the following exact upper bound holds:
\begin{equation}\label{eq:remote-core-factor-exact}
\begin{split}
 \mathbb E|L+Q+gG|+\mathbb E|L-Q-gG|
 \le{}& f(v)-d_3\mathbb E|c|^2-d\mathbb E|U|^2
           -2\sqrt2\,d\,\Re\mathbb E[c\overline C]\\
      &+d_0\mathbb E|Q_0-A|^2-2d\mathbb E|A|^2.
\end{split}
\end{equation}
The coefficients satisfy
\begin{equation}\label{eq:remote-core-factor-ranges}
 \frac d2\le d_3\le2d,\qquad d\le d_0\le4d.
\end{equation}

For an actual $(2,1)$-chaos expansion along $X$, $c$ is constant, the remote
bidegrees of $(C,D,U,A,Q_0)$ are respectively
$(1,1),(2,0),(0,1),(1,0),(2,1)$, and hence
$\mathbb EC=0$ and $\mathbb E[Q_0\overline A]=0$.
In particular, the sufficient condition
\begin{equation}\label{eq:remote-core-factor-cone}
 |c|^2+2\mathbb E|U|^2\ge4\mathbb E|A|^2+8\mathbb E|Q_0|^2
\end{equation}
gives exact domination by the variance-matched independent Gaussian. The
retained deficit is at least
\[
 d\left(\frac{|c|^2}2+\mathbb E|U|^2
       -2\mathbb E|A|^2-4\mathbb E|Q_0|^2\right).
\]
Thus when $A=Q_0=0$, a coherent Wick core may coexist with arbitrary latent
factors $CX+D\overline X+UX^2/\sqrt2$ at every amplitude. Subtracting
$(v+g^2)/\varepsilon$ gives the corresponding exact comparison with the
Gaussian benchmark for every positive price $\varepsilon$.

\paragraph{Proof.}
Write $B=V+g^2$ and $F(u)=2\mathbb E|\sqrt BZ+\sqrt u|$.
The square-root Laplace formula and a Gaussian integral give, for fixed $B=b$,
\begin{equation}\label{eq:remote-core-factor-bernstein}
 F_b(u)=\sqrt{\pi b}+\frac1{\sqrt\pi}
 \int_0^{1/b}(1-e^{-tu})\frac{\sqrt{1-bt}}{t^{3/2}}\,dt.
\end{equation}
At $b=0$ the upper endpoint is infinity. Thus $F$ is increasing and concave,
and $F'(u)\le u^{-1/2}$ for $u>0$. Set $R=|X|^2$, which is exponential with
mean one. Circularity gives the left side of
\eqref{eq:remote-core-factor-exact} as $\mathbb EF(|Q|^2)$, whereas
$\mathbb EF(vR)=f(v)$. The concave tangent at $vR$ therefore gives
\[
 \mathbb EF(|Q|^2)\le f(v)+\mathbb E[F'(vR)|Q|^2]-vf'(v).
\]
The weighted term is integrable by $F'(vR)\le(vR)^{-1/2}$ and the finite
exponential moments of $R$.

Averaging the phase of $X$ leaves only two cross terms, $cH$ with $CX$, and
$Q_0$ with $A(R-1)$. The five diagonal radial weights are
\[
\begin{array}{c|c}
1 &(f(v)-f(0))/v\\
X,\overline X &f'(v)\\
X^2/\sqrt2 &f'(v)+vf''(v)/2\\
R-1 &(f(v)-f(0))/v+vf''(v)\\
H(X) &f'(v)+vf''(v)+v^2f'''(v)/2.
\end{array}
\]
The two off-diagonal weights are $vf''(v)/\sqrt2=-\sqrt2d$ and
$f'(v)-(f(v)-f(0))/v=-d_0$. These identities follow directly from
$\mathbb E[R^j t e^{-tvR}]=j!t/(1+tv)^{j+1}$ applied to the Bernstein
representation \eqref{eq:remote-core-factor-bernstein}. Substitution proves
\eqref{eq:remote-core-factor-exact}.

To prove \eqref{eq:remote-core-factor-ranges}, put
\[
 r=\sqrt{\frac B{B+v}},\qquad
 h=\frac{\sqrt\pi v}{8(B+v)^{3/2}}.
\]
Direct differentiation gives
\[
 d=\mathbb Eh,\qquad
 d_3=\mathbb E\!\left[h\left(\frac12+\frac32r^2\right)\right],\qquad
 d_0=\mathbb E\!\left[\frac{4h}{(1+r)^2}\right].
\]
Since $0\le r<1$, the claimed ranges follow. Under cubic orthogonality, the
exact deficit decomposes as
\begin{align*}
 &d_3|c|^2+d\mathbb E|U|^2-(d_0-2d)\mathbb E|A|^2-d_0\mathbb E|Q_0|^2\\
 &\quad=d\left(\frac{|c|^2}2+\mathbb E|U|^2
                   -2\mathbb E|A|^2-4\mathbb E|Q_0|^2\right)
       +(d_3-d/2)|c|^2+(4d-d_0)(\mathbb E|A|^2+\mathbb E|Q_0|^2),
\end{align*}
whose last two terms are nonnegative. The case $v=0$ is immediate.
\hfill$\square$

The accompanying self-contained script \texttt{check\_core\_fibre.py} checks
51 finite identities in exact arithmetic. The arbitrary-dimensional result
follows from the proof above, not from a pattern enumeration. The shared
visible $A/B$ coefficients cannot be substituted for the latent coefficients:
they depend on the visible variable and invalidate the first averaging step.

\subsubsection{An optimized tangent proves exact K for every remote star kernel}

At the exact peak profile $a=699/1250$, $b=179/1250$ and
$\varepsilon=1891/2500$, let $w_*$ maximize
$\varphi(w)-w/\varepsilon$, where
$\varphi(w)=\sqrt\pi\,\mathbb E\sqrt{V+w}$.
A scalar interval certificate proves
\[
 0.0403448<w_*<0.0403449,
\]
and, at $g=0$, $v=w_*$,
\[
 d>\frac{45}{1000},\quad d_3>\frac{77}{1000},\quad
 2d-d_0>\frac{39}{1000},\quad d_0<\frac{52}{1000}.
\]
For every remote cubic in the expansion above and every independent Gaussian
buffer, put $\theta_g=(1-g^2/w_*)_+$. Then
\begin{equation}\label{eq:remote-star-optimized}
 \boxed{\Lambda(Q+gG)\le S-\frac{\theta_g}{1000}
  \left(77|c|^2+45\mathbb E|U|^2+39\mathbb E|A|^2
                            -52\mathbb E|Q_0|^2\right).}
\end{equation}
In particular, $Q_0=0$ gives exact K for every remote cubic whose Hermite
monomials all involve one chosen remote coordinate, in arbitrary dimension
and at every amplitude. The arbitrary latent coefficients may be dependent.
The centered quadratic fibre $A(|X|^2-1)$ here is remote; it is not the open
visible shared-$A$ sector. More generally,
$77|c|^2+45\mathbb E|U|^2+39\mathbb E|A|^2\ge52\mathbb E|Q_0|^2$
suffices. The factor $\theta_g$ must not be omitted.

For the proof, if $g^2<w_*$ anchor the concave tangent at
$(w_*-g^2)|X|^2$, regardless of the actual variance of $Q$.
Its first-chaos weight is exactly $\varphi'(w_*)=1/\varepsilon$, so the
variance price cancels at every amplitude. In the component formulas above,
$r_g=\sqrt{(V+g^2)/(V+w_*)}$ increases with $g^2$, and the common factor is
$\theta_g$ times its value at $g=0$. Therefore
\[
 d_g=\theta_gd_*,\quad (d_3)_g\ge\theta_g(d_3)_*,\quad
 (d_0)_g\le\theta_g(d_0)_*,\quad
 2d_g-(d_0)_g\ge\theta_g(2d_*-(d_0)_*).
\]
These inequalities prove \eqref{eq:remote-star-optimized}. If $g^2\ge w_*$,
the tangent at zero instead gives, for every square-integrable remote $Q$,
\[
 \Lambda(Q+gG)\le\varphi(g^2)-g^2/\varepsilon
 +(\varphi'(g^2)-1/\varepsilon)\mathbb E|Q|^2\le S.
\]

\paragraph{Scalar certificate.}
The included finite Gaussian-mixture density witness has radial $L^1$ error
at most $e=10^{-11}$ and doubled radius-weighted $L^1$ error at most $e$.
Its finite-cell and infinite-tail replay is included. For the Gaussian-smoothed
twice-norm, the first three derivative kernels have supremum at most
\[
 \frac{\sqrt\pi}{2\sqrt w},\qquad
 \frac{\sqrt{2\pi}}{w^{3/2}},\qquad
 \frac{8\sqrt{2\pi}}{w^{5/2}},
\]
respectively. Indeed the first kernel is the convolution of $1/(2|z|)$ with
the variance-$w$ Gaussian. For $p_w(z)=e^{-|z|^2/w}/(\pi w)$, direct
calculation yields $|\partial_wp_w|\le4p_{2w}/w$ and
$|\partial_w^2p_w|\le32p_{2w}/w^2$, by the elementary bounds
$x+1\le2e^{x/2}$ and $x^2+4x+2\le16e^{x/2}$ for $x\ge0$.
Convolution gives the stated derivative bounds. Multiplication by $e$
therefore bounds every derivative error used by the certificate.

The script \texttt{check\_peak\_anchor.py} encloses the derivative signs at
both rational endpoints of the displayed interval, then certifies all four
rational coefficient bounds throughout that interval. The optimizer is
implicit but rigorously enclosed; the conclusion is exact K, not a slack
version. This result enlarges a remote class and does not prove K for shared
visible $A/B$ couplings or change a numerical Grothendieck endpoint.

\paragraph{Two-coordinate corollary.}
Every remote $q\in H_{2,1}(\mathbb C^2)$ satisfies
\[
 \boxed{\Lambda(q+gG)\le S-\frac{\theta_g}{80}\mathbb E|q|^2.}
\]
To see this, write the six orthonormal coefficients as $c_1,c_2$ for the
diagonal Wick terms, $u_1,u_2$ for
$X^2\overline Y/\sqrt2,Y^2\overline X/\sqrt2$, and $a_1,a_2$ for
$Y(|X|^2-1),X(|Y|^2-1)$. Average \eqref{eq:remote-star-optimized} with $X$
as the chosen coordinate and with $Y$ as the chosen coordinate. The result is
\[
 \Lambda(q+gG)\le S-\frac{\theta_g}{2000}
 \left[25(|c_1|^2+|c_2|^2)+45(|u_1|^2+|u_2|^2)
                         +39(|a_1|^2+|a_2|^2)\right].
\]
The smallest coefficient is $1/80$. Thus this entire class has a global
quantitative proof based on four scalar bounds, rather than a cover of the
cubic coefficient space.

\subsection{Exact complex orbit
constraints}\label{exact-complex-orbit-constraints}

All sphere measures below are normalized. A standard proper complex
Gaussian has density \texttt{exp(-\textbar{}z\textbar{}\textsuperscript{2})/\ensuremath{\pi}}. Gaussian
Hermite projections are denoted by \texttt{P\_\{p,q\}}. The ambient
Gaussian dimension is arbitrary and finite; every inequality extends by
the usual finite-dimensional conditional-expectation approximation.

\subsubsection{1. A six-point complex orbit is
nonvacuous}\label{a-six-point-complex-orbit-is-nonvacuous}

Take the six unit vectors in C\textsuperscript{2} that are eigenvectors of the three Pauli
matrices, two eigenvectors for each matrix. Let \texttt{G} be their Gram
matrix. Then the complex phase optimum of the positive-semidefinite
matrix G is

\[
 \max_{|z_i|=1}z^*Gz=9+3\sqrt6,
\]

whereas its vector SDP optimum is exactly 18. Thus the ratio is
\texttt{2(3-sqrt(6))\ \textgreater{}\ 1}.

\textbf{Proof.} The six vectors form a tight frame with frame operator
\texttt{3I\textsubscript{2}}, so G has eigenvalues 3,3,0,0,0,0. The SDP dual
\texttt{3I\textsubscript{6}-G\ \textgreater{}=\ 0} gives 18; choosing the six original
vectors as the vector assignment gives \texttt{tr(G\textsuperscript{2})=18}.

The scalar optimum is \[
 \left(\max_{\|u\|=1}\sum_i|\langle u,v_i\rangle|\right)^2.
\] For a unit vector u with Bloch vector \texttt{(n\textsubscript{1},n\textsubscript{2},n\textsubscript{3})},
\texttt{sum\ n\_j\textsuperscript{2}=1}, the sum is \[
 \sum_{j=1}^3\sqrt{1+\sqrt{1-n_j^2}}.
\] The function \texttt{x\ -\textgreater{}\ sqrt(1+sqrt(1-x))} is
concave on \texttt{{[}0,1{]}}. Jensen gives at most
\texttt{3\ sqrt(1+sqrt(2/3))}; equality occurs at \texttt{n\_j\textsuperscript{2}=1/3}.
Squaring proves the claim.

This disproves the possibility that all small complex orbit constraints
are vacuous. It is not itself competitive with the existing Grothendieck
lower bound.

\subsubsection{2. The sharp unanchored sphere
hierarchy}\label{the-sharp-unanchored-sphere-hierarchy}

For every disk-valued Gaussian function f and integer
\texttt{m\ \textgreater{}=\ 2}, \[
 \sum_{k\ge0} w_k^{(m)}\|P_{k+1,k}f\|_2^2\le b_m,
 \qquad
 w_k^{(m)}=\frac{m!(k+1)!}{(m+k)!},
\] \[
 b_m=\frac{m\pi}{4}
       \left(\frac{\Gamma(m)}{\Gamma(m+1/2)}\right)^2.
\tag{2.1}
\] Other charge sectors of f are permitted: they simply do not
contribute on the left. In particular, \[
 \sum_{k\ge0}\frac{2}{k+2}\|P_{k+1,k}f\|_2^2\le\frac89.
\tag{2.2}
\] The phase of one Gaussian coordinate gives equality in every member
of (2.1).

\textbf{Proof.} Let \texttt{Z\textsubscript{1},...,Z\_m} be independent copies of the
ambient Gaussian vector. For each realization define
\texttt{F(u)=f(sum\_j\ u\_j\ Z\_j)} on the unit sphere in C\^{}m and \[
 B_j=\int\overline{u_j}F(u)\,du.
\] By duality and rotational invariance, \[
 \|B\|\le\int|u_1|\,du
 =\frac{\Gamma(m)\Gamma(3/2)}{\Gamma(m+1/2)}.
\tag{2.3}
\] For two independent sphere points u,v, put
\texttt{rho=sum\ u\_j\ bar(v\_j)}. Gaussian Hermite orthogonality gives
\[
 \mathbb E\bigl[f(\textstyle\sum u_jZ_j)
       \overline{f(\sum v_jZ_j)}\bigr]
 =\sum_{p,q}\rho^p\bar\rho^q\|P_{p,q}f\|_2^2.
\] Multiplication by \texttt{bar(rho)} and sphere averaging kills all
terms except \texttt{p=q+1}. Since \texttt{\textbar{}rho\textbar{}\textsuperscript{2}} has
Beta(1,m-1) law, \[
 \mathbb E\|B\|^2
 =\sum_{k\ge0}\frac{(k+1)!(m-1)!}{(m+k)!}
                   \|P_{k+1,k}f\|_2^2.
\tag{2.4}
\] Combine (2.3)-(2.4) and multiply by m.

For \texttt{f(z)=phase(z\textsubscript{1})}, almost every orbit is
\texttt{F(u)=phase(sum\ u\_j\ z\_j)}, which attains equality in (2.3);
hence equality also holds after averaging. Equivalently its Hermite
energies are \[
 e_k=\frac\pi4\frac{(1/2)_k^2}{k!(k+1)!},
\] and Gauss summation gives \texttt{sum\ w\_k\^{}(m)e\_k=b\_m}.

\subsubsection{3. A sharp anchored disk
inequality}\label{a-sharp-anchored-disk-inequality}

Let U be uniform on the complex unit disk. For all complex a,b, \[
 \mathbb E|a+bU|\le\sqrt{|a|^2+\frac49|b|^2}.
\tag{3.1}
\] The coefficient \texttt{4/9} is sharp, by taking a=0.

\textbf{Proof.} By scaling and rotation put
\texttt{M(a)=E\textbar{}a+U\textbar{}}, \texttt{a\textgreater{}=0},
\texttt{b=1}. First, \texttt{M(0)=2/3},
\texttt{M\textquotesingle{}(0)=0}, and \[
 M''(a)=\frac1\pi\int_{x^2+y^2\le1}
       \frac{y^2}{((x+a)^2+y^2)^{3/2}}\,dx\,dy\le M''(0)=1.
\] For each fixed y, the integrand as a function of x+a is even and
nonincreasing in absolute value, and its interval of integration is
centered at zero before shifting; this proves the inequality. The
formula can also be justified by first replacing the norm by
\texttt{sqrt(\textbar{}z\textbar{}\textsuperscript{2}+epsilon\textsuperscript{2})} and then taking epsilon
down to zero. Consequently \texttt{M(a)\textless{}=2/3+a\textsuperscript{2}/2}. For
\texttt{0\textless{}=a\textless{}=1}, \[
 a^2+\frac49-\left(\frac23+\frac{a^2}2\right)^2
 =\frac{a^2}3-\frac{a^4}4\ge0.
\]

For \texttt{a\textgreater{}=1}, angular binomial expansion followed by
radial integration gives the absolutely convergent positive-coefficient
series \[
 M(a)=a\sum_{n\ge0}
 \frac{\binom{1/2}{n}^2}{n+1}a^{-2n}.
\] Writing \texttt{d=M(1)-1\textgreater{}=0} gives
\texttt{M(a)\textless{}=a+d/a}. Direct polar integration over the disk
centered at 1 gives \texttt{M(1)=32/(9\ensuremath{\pi})}. Thus \[
 M(a)^2-a^2\le2d+d^2=M(1)^2-1
 <(32/27)^2-1=295/729<4/9,
\] where only \texttt{\ensuremath{\pi}\textgreater{}3} was used. This covers all a.

\subsubsection{4. A conditional constraint retaining the
mean}\label{a-conditional-constraint-retaining-the-mean}

For any measurable \texttt{F:S\textsuperscript{3}\ -\textgreater{}\ C} with
\texttt{\textbar{}F\textbar{}\textless{}=1}, write \[
 m_F=\int F(u)\,du,\qquad B_F=\int\bar u F(u)\,du.
\] Then \[
 |m_F|^2+\frac94\|B_F\|^2\le1.
\tag{4.1}
\]

Indeed the first coordinate of a uniform point on S\textsuperscript{3} is uniform on the
unit disk. The support function of the set of \texttt{(m\_F,B\_F)} is
exactly \texttt{E\textbar{}a+bU\textbar{}}; (3.1) and ellipsoid duality
give (4.1).

For a Gaussian disk-valued function f, apply (4.1) to the orbit from
Section 2 with m=2 and average. The mean term selects charge zero; the B
term selects charge one. This gives \[
 \boxed{\quad
 \sum_{k\ge0}\frac{\|P_{k,k}f\|_2^2}{k+1}
 +\frac94\sum_{k\ge0}\frac{\|P_{k+1,k}f\|_2^2}{k+2}
 \le1.\quad}
\tag{4.2}
\] The first coefficient follows from
\texttt{E\textbar{}rho\textbar{}\^{}(2k)=1/(k+1)}; the second follows
from \texttt{E\textbar{}rho\textbar{}\^{}(2k+2)=1/(k+2)}.

Most importantly, (4.2) can be applied \textbf{conditionally on the
visible coordinate}. The projections then act only in the remote
Gaussian variables. It therefore couples the actual fibre mean to remote
Hermite energies instead of treating their separate bounds as
independently attainable. More generally any independent external
parameter may be held fixed, and any nonnegative measurable weight of
that parameter may be integrated afterwards.

\subsubsection{5. The affine Gaussian moment body, including both
polarizations}\label{the-affine-gaussian-moment-body-including-both-polarizations}

If X is a standard proper Gaussian vector and
\texttt{\textbar{}f(X)\textbar{}\textless{}=1}, let \[
 m=\mathbb Ef,\qquad B=\mathbb E[\bar Xf],\qquad
 D=\mathbb E[Xf].
\] For every affine test \texttt{(a,b,c)}, the exact support function of
the attainable moment set is \[
 \sup_{|f|\le1}\Re\mathbb E[\overline{a+b^TX+c^T\bar X}\,f]
 =\mathbb E|a+b^TX+c^T\bar X|.
\tag{5.1}
\] Thus an exact finite-dimensional convex moment constraint is
available without a price relaxation. Its dependence is only on the mean
a, variance
\texttt{v=\textbar{}\textbar{}b\textbar{}\textbar{}\textsuperscript{2}+\textbar{}\textbar{}c\textbar{}\textbar{}\textsuperscript{2}},
and pseudocovariance \texttt{p=2b\^{}Tc} (including the relative phase
of a\textsuperscript{2} and p).

For one polarization there is the simpler sharp bound \[
 \mathbb E|a+b^TX|\le\sqrt{|a|^2+\frac\pi4\|b\|^2},
 \qquad |m|^2+\frac4\pi\|B\|^2\le1.
\tag{5.2}
\]

\textbf{Proof of the support bound.} For standard G put
\texttt{M(a)=E\textbar{}a+G\textbar{}}, with a real and nonnegative. In
polar coordinates about the shifted origin, conditionally on radius R=r
the angle has density proportional to \texttt{exp(2ar\ cos(theta))}. Its
conditional mean cosine is nondecreasing in r because its derivative is
\texttt{2a\ Var(cos(theta)\textbar{}R=r)\textgreater{}=0}. Hence \[
 a=\mathbb E[R\cos\theta]\ge
        \mathbb ER\,\mathbb E\cos\theta=M(a)M'(a).
\] Integration from zero gives
\texttt{M(a)\textsuperscript{2}\textless{}=a\textsuperscript{2}+M(0)\textsuperscript{2}=a\textsuperscript{2}+\ensuremath{\pi}/4}. Scaling proves (5.2);
ellipsoid duality proves the moment statement.

For both polarizations a valid explicit bound retaining pseudocovariance
is \[
 \boxed{\quad
 \bigl(\mathbb E|a+b^TX+c^T\bar X|\bigr)^2
 \le |a|^2+\frac\pi4v+\left(1-\frac\pi4\right)|p|.
 \quad}
\tag{5.3}
\] To prove this, decompose the centered scalar Gaussian as the sum of
an independent proper Gaussian of variance
\texttt{v-\textbar{}p\textbar{}} and a rank-one real Gaussian of
variance \texttt{\textbar{}p\textbar{}}. Apply (5.2) conditionally on
the real Gaussian and then apply Jensen to the square root. This gives
precisely (5.3). The exact support (5.1) can be used instead when the
extra loss is material.

There is no valid uniform extension assigning the improved coefficient
to both B and D separately. On S\textsuperscript{3} take
\texttt{F\_delta(u)=phase(1+i\ delta\ Re(u\textsubscript{1}))}. As delta tends to zero,
\[
 |\mathbb EF_\delta|^2=1-\delta^2/4+O(\delta^4),
 \quad \|B\|^2+\|D\|^2=\delta^2/8+O(\delta^4).
\] Thus
\texttt{\textbar{}m\textbar{}\textsuperscript{2}+c(\textbar{}\textbar{}B\textbar{}\textbar{}\textsuperscript{2}+\textbar{}\textbar{}D\textbar{}\textbar{}\textsuperscript{2})\textless{}=1}
requires \texttt{c\textless{}=2}, exactly the ordinary Bessel
coefficient. For a proper Gaussian the analogous calculation requires
\texttt{c\textless{}=1}. This is why the pseudocovariance in (5.1) or
(5.3) cannot be dropped.

\paragraph{An explicit inequality on both first-chaos moment
vectors}\label{an-explicit-inequality-on-both-first-chaos-moment-vectors}

The support bound (5.3) has a useful explicit dual. Put \[
 V=\|B\|^2+\|D\|^2,\qquad P=2|B^TD|,
 \quad s_1^2=(V+P)/2,\quad s_2^2=(V-P)/2,
 \quad \beta=\pi/2-1.
\] Then \[
 \boxed{\quad |m|^2+\Phi(s_1,s_2)\le1,\qquad
 \Phi(s_1,s_2)=
 \begin{cases}
  s_1^2+s_2^2/\beta,&s_2\le\beta s_1,\\
  (s_1+s_2)^2/(1+\beta),&s_2\ge\beta s_1.
 \end{cases}\quad}
\tag{5.4}
\] Both formulas agree at the boundary. In particular, when
\texttt{B\^{}TD=0} this gives \[
 |m|^2+\frac4\pi(\|B\|^2+\|D\|^2)\le1.
\tag{5.5}
\] At the opposite, rank-one extreme \texttt{P=V}, it gives ordinary
Bessel. The coefficient \texttt{Phi/V} lies between 1 and \texttt{4/\ensuremath{\pi}}
and depends on the actual output polarization rather than being assigned
separately to the two charges.

\textbf{Proof.} Write the complex Gaussian vector as
\texttt{(N\textsubscript{1}+iN\textsubscript{2})/sqrt(2)}, with standard real Gaussian coordinates N,
and let \[
 T=\mathbb E[(\Re f,\Im f)^TN^T].
\] Its singular values are precisely \texttt{s\textsubscript{1},s\textsubscript{2}}, since
\texttt{\textbar{}\textbar{}T\textbar{}\textbar{}\_F\textsuperscript{2}=V} and the
difference of the eigenvalues of \texttt{TT\^{}T} is P. For an arbitrary
real \texttt{2\ by\ 2n} matrix A, the centered complex Gaussian
represented by AN has covariance eigenvalues
\texttt{sigma\textsubscript{1}(A)\textsuperscript{2},sigma\textsubscript{2}(A)\textsuperscript{2}}. Thus (5.3) reads \[
 \mathbb E|a+AN|\le
 \sqrt{|a|^2+g(A)^2},\qquad
 g(A)^2=\sigma_1(A)^2+\beta\sigma_2(A)^2.
\] Here g is a norm: its square is
\texttt{beta\ \textbar{}\textbar{}A\textbar{}\textbar{}\_F\textsuperscript{2}+(1-beta)\textbar{}\textbar{}A\textbar{}\textbar{}\_op\textsuperscript{2}}.
Ellipsoid duality in the mean coordinate, followed by norm duality in A,
gives \texttt{\textbar{}m\textbar{}\textsuperscript{2}+g\_*(T)\textsuperscript{2}\textless{}=1}.
Singular-value alignment reduces the dual norm calculation to \[
 g_*(T)=\max\{s_1x+s_2y:x\ge y\ge0,\ x^2+\beta y^2\le1\}.
\] When \texttt{s\textsubscript{2}\textless{}=beta\ s\textsubscript{1}}, the unconstrained ellipsoid
maximizer satisfies \texttt{y\textless{}=x} and gives the first branch
of (5.4). Otherwise the maximum lies on \texttt{x=y}, giving the second
branch. This proves the formula.

All of (5.1)-(5.5) hold pointwise after conditioning on an independent
visible variable. They can therefore be imposed on its actual
conditional mean and conditional first-chaos coefficient vectors, and
subsequently integrated with any nonnegative weight. They do not assert
anything about arbitrary cubics or their laws.

\subsubsection{6. Use and limitation}\label{use-and-limitation}

The new inequalities constrain bounded maximizing responses directly,
before a free cubic variable is introduced. They are therefore
legitimate possible replacements for some of the stronger assertions in
K. For \texttt{h=(f+g)/2}, \texttt{k=(f-g)/2}, apply an energy
inequality to both f and g and average; the parallelogram identity
converts it to the sum of the h and k energies.

A global positive Hermite multiplier from (2.1) alone is not enough: the
current diffuse response has slack in the corresponding constraints, and
the parent task's probe found the resulting simple shifts worsen the row
objective. The promising interface is (4.2) conditionally on the visible
coordinate, or (5.1) with the visible fibre mean retained. Turning
either into a smaller certified row cap still requires an optimization
and a complete envelope. No numerical endpoint follows from this note
alone.
